\chapter{Heel Pain}
\index{Heel Pain}

\section*{Definition and Overview}
Heel pain is pain felt under or behind the heel, most commonly caused by plantar fasciitis, an inflammation of the thick band of tissue that runs along the sole of the foot. It can also arise from Achilles tendon problems, bursitis, stress fractures, and nerve entrapment. Heel pain affects people of all ages and is particularly common in runners, people who stand for long periods, and those who are overweight. Most cases respond well to simple treatment, including stretching, footwear changes, and activity modification, though recovery often takes months.

\section*{Epidemiology}
Plantar fasciitis is one of the most common causes of foot pain, affecting around one in ten people at some point in life. It is most common in people aged 40--60 and is strongly associated with prolonged standing, obesity, running, and occupations that involve long periods on the feet. Achilles tendon disorders also become more common with age and in active people. Heel pain accounts for a large number of GP, podiatry, and physiotherapy consultations, and many people self-manage without seeking care.

\section*{Aetiology and Pathophysiology}
\begin{itemize}
    \item \textbf{Plantar fasciitis:} repetitive strain and micro-tears where the plantar fascia attaches to the heel bone, aggravated by overuse, poor foot mechanics, tight calf muscles, and inappropriate footwear
    \item \textbf{Achilles tendinopathy:} overload and degeneration of the Achilles tendon, causing pain at the back of the heel
    \item \textbf{Retrocalcaneal bursitis:} inflammation of the bursa behind the heel
    \item \textbf{Stress fracture:} overload of the heel bone, especially in runners
    \item \textbf{Nerve entrapment:} compression of nerves around the heel
    \item \textbf{Systemic conditions:} inflammatory arthritis, gout, and diabetes-related foot problems
    \item \textbf{Pathophysiology:} repeated loading of the foot structures causes micro-injury, inflammation, and degeneration that outpace repair
\end{itemize}

\section*{Clinical Features}
\begin{itemize}
    \item Sharp pain under the heel with the first steps in the morning, easing after a few minutes
    \item Pain after prolonged standing or sitting, worse with weight-bearing
    \item Pain that worsens with running, walking on hard surfaces, or barefoot activity
    \item Pain at the back of the heel with Achilles problems
    \item Tenderness to pressure over the heel
    \item Stiffness of the foot and calf
    \item Symptoms typically develop gradually over weeks to months
\end{itemize}

\section*{Differential Diagnosis}
\begin{itemize}
    \item Plantar fasciitis is the commonest cause of inferior heel pain
    \item Achilles tendinopathy and bursitis cause posterior heel pain
    \item Stress fractures cause bone pain with activity
    \item Fat pad atrophy causes pain under the heel in older people
    \item Nerve entrapment causes burning, tingling, or radiating pain
    \item Inflammatory arthritis, gout, and infection must be considered, especially with swelling, redness, or fever
\end{itemize}

\section*{When to Seek Medical Care}
See a doctor, podiatrist, or physiotherapist if heel pain persists beyond a few weeks, interferes with daily activities, or does not respond to simple self-care. Seek urgent review for pain after a fall, severe swelling, or an inability to bear weight.

\textbf{Contact your local emergency services for:}
\begin{itemize}
    \item Severe heel or foot pain with inability to walk
    \item Sudden "snap" in the Achilles area with inability to push off the foot
    \item A hot, red, swollen heel with fever, which may indicate infection
    \item Numbness, tingling, or loss of sensation in the foot
    \item Injury from a fall with severe pain or deformity
\end{itemize}

\section*{Diagnosis and Investigations}
\begin{itemize}
    \item \textbf{Clinical examination:} location of pain, tenderness, foot mechanics, and range of movement usually make the diagnosis
    \item \textbf{History:} activity levels, footwear, occupation, and duration of symptoms
    \item \textbf{Imaging:} X-ray to exclude fracture and bone problems; ultrasound and MRI for soft tissue and tendon assessment
    \item \textbf{Ultrasound:} can visualise the plantar fascia and Achilles tendon
    \item \textbf{Blood tests:} for inflammatory conditions and gout when suspected
    \item \textbf{Review:} gait and footwear assessment to identify contributing factors
\end{itemize}

\section*{Management}
\begin{itemize}
    \item \textbf{Rest and activity modification:} reduce aggravating activities and avoid barefoot walking
    \item \textbf{Stretching:} calf and plantar fascia stretches, particularly before getting out of bed
    \item \textbf{Footwear:} supportive shoes with cushioning, and orthotics or insoles where helpful
    \item \textbf{Ice:} ice massage reduces pain and inflammation
    \item \textbf{Pain relief:} simple analgesics and short-term anti-inflammatories
    \item \textbf{Physiotherapy:} strengthening, loading programs, and manual therapy
    \item \textbf{Shockwave therapy:} effective for chronic plantar fasciitis in some patients
    \item \textbf{Injections:} corticosteroid injections for persistent cases, used judiciously
    \item \textbf{Surgery:} rarely needed for plantar fascia release or Achilles repair
\end{itemize}

\section*{Complications}
\begin{itemize}
    \item Chronic pain lasting a year or more in a minority of people
    \item Altered gait, leading to knee, hip, and back problems
    \item Reduced activity, weight gain, and deconditioning
    \item Rupture of the Achilles tendon with sudden, severe injury
    \item Missed stress fractures or inflammatory conditions
    \item Work disability in occupations requiring prolonged standing
\end{itemize}

\section*{Prognosis and Outcomes}
Heel pain resolves in the majority of people with consistent treatment, although recovery is often slow, typically taking 3--12 months for plantar fasciitis. Most people improve with stretching, footwear modification, and load management alone, and physiotherapy-led programs have high success rates. Surgery is rarely required. Early treatment, correction of contributing factors, and patience are the keys to a good outcome; untreated heel pain tends to persist and can become chronic.

\section*{Prevention and Self-Care}
\begin{itemize}
    \item Wear supportive, well-cushioned shoes and replace worn footwear
    \item Stretch the calves and plantar fascia regularly, especially before activity
    \item Maintain a healthy weight to reduce load on the feet
    \item Build up running and exercise volume gradually
    \item Avoid prolonged standing on hard surfaces and take regular breaks
    \item Warm up before sport and cool down afterwards
    \item Seek early treatment rather than waiting for pain to become chronic
\end{itemize}

\section*{Sources and Further Reading}
The following reputable sources provide further information for healthcare students and patients seeking reliable, current advice about heel pain.

\begin{itemize}
    \item Healthdirect. \textit{Heel pain: causes and treatment.} https://www.healthdirect.gov.au/heel-pain
    \item Sports Medicine Australia. \textit{Plantar fasciitis and foot injuries.} https://sma.org.au/
    \item Australian Physiotherapy Association. \textit{Heel pain information.} https://australian.physio/
    \item NHS. \textit{Heel pain: causes and treatment.} https://www.nhs.uk/conditions/heel-pain/
    \item Mayo Clinic. \textit{Plantar fasciitis: symptoms and causes.} https://www.mayoclinic.org/
    \item MSD Manual. \textit{Plantar fasciitis.} https://www.msdmanuals.com/
\end{itemize}
